% GENERATED FILE. Edit canonical lessons, then run npm run generate:books.
% canonical-source-hash: fnv1a64:36ddd1f407ad8c49
% canonical-lessons: UR-C59-almari, UR-C59-farsh, UR-C59-sirhi, UR-C59-aina, UR-C59-batti

\chapter{Cupboard, Floor, Stairs, Mirror, Light}
\label{ch:ur-cupboard-floor-stairs-mirror-light}
\hlchaptermodality{59}

\begin{chapteropening}
\textbf{By the end of this chapter:} \emph{I can say a cupboard, the floor, stairs, a ladder, a mirror and a light, a lamp.}

Complete the last lesson of chapter 59: I can say a cupboard, the floor, stairs, a ladder, a mirror and a light, a lamp.
\end{chapteropening}

\section[\texorpdfstring{\ur{الماری} (almārī)}{الماری (almārī)}]{\ur{الماری} (almārī) — a cupboard}

\label{lesson:UR-C59-almari}



\begin{quote}
Before the new one: say the Urdu for a roof; a ceiling, then the Urdu for a bed.
\end{quote}

\subsection*{You'll want to know: \ur{الماری}}
\textbf{\ur{الماری}} — \emph{almārī} — \textquotedblleft{}a cupboard\textquotedblright{}.

\begin{cousinweb}
\textbf{\ur{الماری}} (\emph{almārī}) is usually traced to Portuguese \emph{armário}, \textquotedblleft{}cupboard\textquotedblright{}.
\end{cousinweb}

\begin{grammarlens}[title={What you've built}]
One more word for this chapter.
\end{grammarlens}

\subsection*{Your turn}
\begin{itemize}
\raggedright
  \item \emph{Say it:} \emph{almārī}
  \item \emph{Say it:} \emph{almārī}, once more
  \item \emph{Say it:} say \emph{almārī}
  \item \emph{Recall:} say the Urdu for a roof; a ceiling, then the Urdu for a bed, then say \emph{almārī} again
\end{itemize}

\begin{tcolorbox}[breakable,colback=teal!4,colframe=teal!35!black,title={Before you move on}]
What does \ur{الماری} mean? (\textquotedblleft{}A cupboard\textquotedblright{}.) Say it once more.
\end{tcolorbox}
\section[\texorpdfstring{\ur{فرش} (farsh)}{فرش (farsh)}]{\ur{فرش} (farsh) — the floor}

\label{lesson:UR-C59-farsh}



\begin{quote}
Before the new one: say the Urdu for a bed, then the Urdu for a cupboard.
\end{quote}

\subsection*{You'll want to know: \ur{فرش}}
\textbf{\ur{فرش}} — \emph{farsh} — \textquotedblleft{}the floor\textquotedblright{}.

\begin{cousinweb}
\textbf{\ur{فرش}} (\emph{farsh}) is a loan from Arabic \emph{farsh}, \textquotedblleft{}anything spread, a carpet\textquotedblright{}, which reached Urdu through Persian. Its Arabic root is \emph{f-r-sh} (Lane). In Urdu it came to mean \textquotedblleft{}the floor\textquotedblright{}.
\end{cousinweb}

\begin{grammarlens}[title={What you've built}]
One more word for this chapter.
\end{grammarlens}

\subsection*{Your turn}
\begin{itemize}
\raggedright
  \item \emph{Say it:} \emph{farsh}
  \item \emph{Say it:} \emph{farsh}, once more
  \item \emph{Say it:} say \emph{farsh}
  \item \emph{Recall:} say the Urdu for a bed, then the Urdu for a cupboard, then say \emph{farsh} again
\end{itemize}

\begin{tcolorbox}[breakable,colback=teal!4,colframe=teal!35!black,title={Before you move on}]
What does \ur{فرش} mean? (\textquotedblleft{}The floor\textquotedblright{}.) Say it once more.
\end{tcolorbox}
\section[\texorpdfstring{\ur{سیڑھی} (sīṛhī)}{سیڑھی (sīṛhī)}]{\ur{سیڑھی} (sīṛhī) — stairs, a ladder}

\label{lesson:UR-C59-sirhi}



\begin{quote}
Before the new one: say the Urdu for a cupboard, then the Urdu for the floor.
\end{quote}

\subsection*{You'll want to know: \ur{سیڑھی}}
\textbf{\ur{سیڑھی}} — \emph{sīṛhī} — \textquotedblleft{}stairs, a ladder\textquotedblright{}.

\begin{cousinweb}
\textbf{\ur{سیڑھی}} (\emph{sīṛhī}) is a tadbhava, a word inherited and reshaped by sound change. It comes from \emph{*śrīḍhi}, \textquotedblleft{}resting-place, ladder\textquotedblright{}, a form CDIAL reconstructs from its descendants (CDIAL 12709).
\end{cousinweb}

\begin{grammarlens}[title={What you've built}]
One more word for this chapter.
\end{grammarlens}

\subsection*{Your turn}
\begin{itemize}
\raggedright
  \item \emph{Say it:} \emph{sīṛhī}
  \item \emph{Say it:} \emph{sīṛhī}, once more
  \item \emph{Say it:} say \emph{sīṛhī}
  \item \emph{Recall:} say the Urdu for a cupboard, then the Urdu for the floor, then say \emph{sīṛhī} again
\end{itemize}

\begin{tcolorbox}[breakable,colback=teal!4,colframe=teal!35!black,title={Before you move on}]
What does \ur{سیڑھی} mean? (\textquotedblleft{}Stairs, a ladder\textquotedblright{}.) Say it once more.
\end{tcolorbox}
\section[\texorpdfstring{\ur{آئینہ} (āīnā)}{آئینہ (āīnā)}]{\ur{آئینہ} (āīnā) — a mirror}

\label{lesson:UR-C59-aina}



\begin{quote}
Before the new one: say the Urdu for the floor, then the Urdu for stairs.
\end{quote}

\subsection*{You'll want to know: \ur{آئینہ}}
\textbf{\ur{آئینہ}} — \emph{āīnā} — \textquotedblleft{}a mirror\textquotedblright{}.

\begin{cousinweb}
\textbf{\ur{آئینہ}} (\emph{āīnā}) is a loan, usually traced to Persian \emph{āīna}, \textquotedblleft{}mirror\textquotedblright{}.
\end{cousinweb}

\begin{grammarlens}[title={What you've built}]
One more word for this chapter.
\end{grammarlens}

\subsection*{Your turn}
\begin{itemize}
\raggedright
  \item \emph{Say it:} \emph{āīnā}
  \item \emph{Say it:} \emph{āīnā}, once more
  \item \emph{Say it:} say \emph{āīnā}
  \item \emph{Recall:} say the Urdu for the floor, then the Urdu for stairs, then say \emph{āīnā} again
\end{itemize}

\begin{tcolorbox}[breakable,colback=teal!4,colframe=teal!35!black,title={Before you move on}]
What does \ur{آئینہ} mean? (\textquotedblleft{}A mirror\textquotedblright{}.) Say it once more.
\end{tcolorbox}
\section[\texorpdfstring{\ur{بتی} (battī)}{بتی (battī)}]{\ur{بتی} (battī) — a light, a lamp}

\label{lesson:UR-C59-batti}



\begin{quote}
Before the new one: say the Urdu for stairs, then the Urdu for a mirror.
\end{quote}

\subsection*{You'll want to know: \ur{بتی}}
\textbf{\ur{بتی}} — \emph{battī} — \textquotedblleft{}a light, a lamp\textquotedblright{}.

\begin{cousinweb}
\textbf{\ur{بتی}} (\emph{battī}) is a tadbhava, a word inherited from Sanskrit and reshaped by sound change. It comes from \emph{varti}, \textquotedblleft{}wick; lamp\textquotedblright{} (CDIAL 11359).
\end{cousinweb}

\begin{grammarlens}[title={What you've built}]
One more word for this chapter.
\end{grammarlens}

\subsection*{Your turn}
\begin{itemize}
\raggedright
  \item \emph{Say it:} \emph{battī}
  \item \emph{Say it:} \emph{battī}, once more
  \item \emph{Say it:} say \emph{battī}
  \item \emph{Recall:} say the Urdu for stairs, then the Urdu for a mirror, then say \emph{battī} again
\end{itemize}

\begin{tcolorbox}[breakable,colback=teal!4,colframe=teal!35!black,title={Before you move on}]
What does \ur{بتی} mean? (\textquotedblleft{}A light, a lamp\textquotedblright{}.) Say it once more.
\end{tcolorbox}
